\documentclass[12pt,a4paper]{article}

\usepackage[ngerman]{babel}
\usepackage[T1]{fontenc}
\usepackage{graphicx}
\usepackage{amsmath}
\usepackage{float}
\usepackage{csquotes}
\usepackage[backend=biber,style=numeric]{biblatex}
\usepackage{tikz}

\addbibresource{literatur.bib}
\title{Aufbau einer Mini-Dokumentation}
\author{Muhammed Hamzic}
\date{\today}
\begin{document}
\maketitle
\section{Einleitung}
Erstellen einer Mini-Dokumentation.
\section{Grundlagen}    
\subsection{IP-Adressen}
Jedes Gerät in einem Netzwerk benötigt
eine IP-Adresse, damit es erreichbar ist.
Das Internet Protocol wurde im RFC 791
beschrieben \cite{postel1981}.
\subsection{Router und Switch}
Ein Switch verbindet Geräte innerhalb
eines lokalen Netzwerks. Ein Router
verbindet verschiedene Netzwerke miteinander.


\section{Hardware}
Wichtige Hardware in unserem Projekt xy. 

\begin{itemize}
    \item Router
    \item Switch
    \item Raspberry Pi
\end{itemize}


\section{Netzwerkaufbau}

In Abbildung~\ref{fig:netzwerk} sieht man
einen unseren Netzwerkaufbau.

\begin{figure}[H]
\centering
\begin{tikzpicture}
\node[draw, rectangle] (router) at (0,2) {Router};
\node[draw, rectangle] (switch) at (0,0) {Switch};
\node[draw, rectangle] (raspberry) at (4,-2) {Raspberry Pi};
\node[draw, rectangle] (raspberry2) at (-4,-2) {Raspberry Pi 2};

\draw (router) -- (switch);
\draw (switch) -- (raspberry);
\draw (switch) -- (raspberry2);
\end{tikzpicture}
\caption{Einfacher Netzwerkaufbau}
\label{fig:netzwerk}
\end{figure}



\section{Vergleich der Netzwerkgeräte}

Die folgende Tabelle zeigt verschiedene
Netzwerkgeräte und ihre Aufgaben.

\begin{table}[H]
\centering
\begin{tabular}{|l|l|}
\hline
Gerät & Aufgabe \\ \hline
Router & Netzwerke verbinden \\ \hline
Switch & Geräte verbinden \\ \hline
Raspberry Pi & Mikrocontroller \\ \hline
\end{tabular}
\caption{Vergleich von Netzwerkgeräten}
\label{tab:geraete}
\end{table}

\section{Berechnung der Übertragungszeit}

Die Übertragungszeit kann mit
Gleichung~\ref{eq:zeit} berechnet werden.

\begin{equation}
t = \frac{D}{R}
\label{eq:zeit}
\end{equation}

Dabei ist $t$ die Zeit in Sekunden,
$D$ die Datenmenge in Bit und
$R$ die Datenrate in Bit pro Sekunde.


\printbibliography

\end{document}

